\section{材料与方法}

\subsection{范围与证据模型}

TE-KG 以人类 TE 为共享标识符，连接原本彼此独立的证据层。文献层保存从论文中提取的生物医学实体和关系；分类层表示 TE 记录之间的分类关系；序列和基因组记录提供来自不同上游来源的参考性或代表性信息；表达矩阵保存样本情境下的测量值；共表达网络则概括 TE 与基因特征在特定情境中的秩相关。系统并未为这些证据层赋予相同的证据含义，而是由各运行时入口标明来源类型，并限制描述结果时可以使用的语言。

\subsection{文献收集、信息提取与标准化}

项目使用已记录的 TE 名称、别名和 TE 相关术语组合检索 PubMed 中的人类 TE 文献。归档流程覆盖截至 2026 年 4 月 13 日可获得的论文。候选记录根据 TE 白名单筛选并在纳入前进行标准化。最终保留的语料库包含 2,308 篇论文，与当前 Neo4j 数据库快照中的 Paper 节点数一致。

保留的文本和元数据通过带约束的 DeepSeek-V3 提示词处理，用于识别以 TE 为中心的生物医学陈述及其参与实体。提取实体被映射到 11 类生物医学实体：碳水化合物、疾病、功能、基因、脂质、突变、肽、药物、蛋白质、RNA 和毒素。每条关系记录保留标准化谓词、一个或多个支持该关系的 PubMed 标识符（PMID），以及支持 PMID 的数量。自动标准化包括使用 80\% 的 Ratcliff--Obershelp 相似度阈值进行名称匹配；当自动解析不足时，再进行人工映射。图导入前还通过人工整理审核保留的实体和关系。

归档方法记录记载了三轮人工审核，每轮包含 50 条记录。本文草稿不将这些百分比作为定量性能结论，因为目前尚未一致恢复抽样标识符、抽样种子、准确分母和中间筛选计算。\draftnote{锁定本小节前，需恢复文献检索式、准确的 DeepSeek-V3 端点或模型版本、中间记录数以及抽样审核决定。}

\subsection{TE 分类、序列与基因组记录}

TE 身份和分类记录来自一个 RepBase 衍生的本地来源，基因组位置则来自 RepeatMasker 衍生的 RMSK 注释来源。Repbase 是重复序列家族和共识序列的经典参考资源 \citep{Bao2015Repbase,Kojima2018HumanRepbase}；Dfam 提供互补的家族模型框架 \citep{Wheeler2013Dfam}。TE-KG 保留这些来源角色之间的区别。RepBase 衍生序列被描述为共识序列或参考序列，RMSK 衍生位置被描述为代表性基因组记录。两者都不被表述为活跃、具有多态性或样本特异性插入的完整清单。

分类体系由 Neo4j 支持的 API 提供，而不是使用第二套浏览器端事实来源。分类边使用明确的图关系 \texttt{CLASSIFIED\_AS}、\texttt{SUBFAMILY\_OF} 和 \texttt{HAS\_SUBCATEGORY}。分类 API 可以提供所有受支持 TE 记录的当前树，也可以提供网页界面使用的 RMSK 加 RepBase 视图。\draftnote{补充准确的 RepeatMasker 版本或构建、RepBase 版本、获取日期以及许可和再分发条款。}

\subsection{表达数据集与情境定义}

表达层包含三个当前使用的矩阵。正常组织矩阵含 37,868 个特征和 205 个样本；正常原代细胞矩阵含 37,868 个特征和 307 个样本；癌细胞系矩阵含 37,868 个特征和 646 个样本。三个矩阵共计 1,158 个样本。特征标识符同时包括 TE 或重复序列特征和基因，因此这些矩阵不会被描述为仅包含基因的表达数据。

正常组织数据与 E-MTAB-1733 和 E-MTAB-2836 研究相关 \citep{Edqvist2015SkinTranscriptome,Uhlen2015HumanProteome}。癌细胞系数据与 PRJNA523380 和 Cancer Cell Line Encyclopedia 相关 \citep{Ghandi2019CCLE}。当前正常原代细胞矩阵与 SRP013565 相关；该 ENCODE RNA-seq 研究包含异质性较高的实验。\draftnote{冻结三个矩阵的“登录号到本地样本”对应清单，并确定 SRP013565 的准确子集及主要引用。}

元数据审核发现，正常组织元数据中存在重复的运行标识符，并且有五个矩阵运行记录无法在元数据中匹配：ERR579126、ERR579137、ERR579144、ERR579145 和 ERR579154。这些问题仍作为限制保留，必须在公开发布清单中解决或明确排除。运行时键 \texttt{normal\_cell\_line} 实际指正常原代细胞情境；论文中使用“正常原代细胞”，以避免与转化细胞系混淆。

\subsection{特定情境下的共表达网络}

共表达分别在正常组织、正常原代细胞和癌细胞系情境中计算。丰度值按 $\log_{2}(\mathrm{count}+1)$ 转换，TE-基因关联使用 Spearman 秩相关衡量。候选边在 $|r| \geq 0.4$ 且 Benjamini--Hochberg 错误发现率（FDR）不高于 0.05 时保留 \citep{Benjamini1995FDR}。保留的正相关边用于 Louvain 社区发现，随机种子为 42，分辨率为 1.8 \citep{Blondel2008Louvain}。系统将这些边解释为统计相关，而不是 TE 调控基因的证据。

完整离线网络输出作为来源记录和导入输入保留。正式界面提供从这些结果中筛选并经批准的、由 MySQL 支持的展示网络。当前展示约定将单次返回的网络限制为最多 50 个节点和 150 条边，以便保持交互响应，同时避免将所展示子图描述为完整相关网络。本文报告的目录版本为 \texttt{v1\_abs0.4\_fdr0.05\_res1.8}。

\subsection{图数据库、关系型存储与网页应用}

正式知识图谱存储在 Neo4j 数据库 \texttt{tekg3} 中。MySQL 存储版本化的 Browse 目录、表达数据和经批准的共表达数据集。PHP API 负责访问两种数据库系统，浏览器端 JavaScript 则负责渲染表格、图表、分类视图以及交互式 G6 或 Canvas 图。该架构避免把大型表达矩阵和展示网络记录强行放入文献图谱，同时保留统一的 TE 中心化界面。

知识图谱界面支持实体搜索、关系筛选、交互式节点扩展、关系图例和 PMID 证据链接。Path 界面用于搜索所选实体之间的连接。Browse、分类体系和 Expression 提供 TE 身份与定量表达谱的结构化访问，Download 则提供可复用的图、分类和表达文件。论文中所有界面相关计数均来自带日期的只读数据库查询或实时 API 响应，而不是直接抄录界面标签。

\subsection{受证据边界约束的自然语言检索}

Agent 和 DeepThink 提供两种访问本地资源的自然语言模式。二者使用包含 12 类证据检索角色的目录，覆盖实体解析、图、路径、序列、分类、基因组、表达、共表达和文献任务。语言模型根据问题选择证据，服务器则执行路由限制、证据充分性检查和受限回退。回答可以包含 PubMed 链接，也可以使用当前浏览器对话中的前文。对话上下文不会在刷新页面、打开新标签页或启动新会话后保留。

智能界面被视为访问层，而不是主要科学结果。内部路由标签、插件名称和开发诊断信息会从面向用户的文字中移除。因此，本文仅报告已实现的行为与测试覆盖，不给出单一的总体准确率。已知的局部限制仍明确披露，包括部分情况下只能进行标题层级的文献检索，以及带疾病限定条件的检索仍可能不完整。

\subsection{验证与快照策略}

数据库计数于 2026 年 7 月 31 日通过只读 Neo4j 查询、实时 PHP API 和仓库检查进行验证。每条生物学关系按其存储方向计数一次。另一种无向遍历计数会将每条已存储的有向关系计算两次，因此未被采用。API 契约、运行时数据库配置、图渲染、分类事实来源、共表达目录一致性以及 Agent/DeepThink 行为均通过仓库中的自动化与浏览器回归脚本检查。与环境相关的响应时间未被用作生物学或性能结论。
